% !TeX program = xelatex
\documentclass[a4paper,11pt]{article}

% --------------------------------------------------
% Engine requirement
% --------------------------------------------------
\usepackage{iftex}
\ifPDFTeX
  \errmessage{This manual must be compiled with XeLaTeX or LuaLaTeX}
\fi

% --------------------------------------------------
% Packages
% --------------------------------------------------
\usepackage{fontspec}
\usepackage{geometry}
\usepackage{xcolor}
\usepackage{hyperref}
\usepackage{booktabs}
\usepackage{enumitem}
\usepackage{fancyvrb}
\usepackage[font=noto, language=burmese]{../myanmar-script}

% --------------------------------------------------
% Page layout
% --------------------------------------------------
\geometry{margin=25mm}
\setlist{itemsep=2pt, topsep=4pt}
\emergencystretch=3em

% --------------------------------------------------
% Fonts for the manual
% --------------------------------------------------
\setmainfont{Latin Modern Roman}
\setsansfont{Latin Modern Sans}
\setmonofont{Latin Modern Mono}

% --------------------------------------------------
% Hyperref setup
% --------------------------------------------------
\hypersetup{
  colorlinks=true,
  linkcolor=black,
  urlcolor=blue,
  pdftitle={myanmar-script Package Manual},
  pdfauthor={Myanmar TeX Project}
}

% --------------------------------------------------
% Verbatim setup
% --------------------------------------------------
\fvset{
  fontsize=\small,
  frame=single
}

% --------------------------------------------------
% Small helpers
% --------------------------------------------------
\newcommand{\pkg}[1]{\textsf{#1}}
\newcommand{\opt}[1]{\texttt{#1}}
\newcommand{\cmd}[1]{\texttt{\textbackslash #1}}
\newenvironment{burmese}
  {\par\smallskip\begin{myanmar}\sloppy}
  {\par\end{myanmar}\fussy\smallskip}

% --------------------------------------------------
% Title info
% --------------------------------------------------
\title{\textbf{The \pkg{myanmar-script} Package}\\
Small helpers for Myanmar-script (Burmese-focused)\\[4pt]
\large\textmyanmar{LaTeX အတွက် မြန်မာစာလုံးရေးစနစ် အထောက်အပံ့}}
\author{Myanmar TeX Project}
\date{Version 1.0 --- 2026/05/29}

\begin{document}

\maketitle

% ==================================================
\section{Introduction}

The \pkg{myanmar-script} package provides a compact interface for typesetting
Myanmar-script (Burmese-focused) text in LaTeX documents. It is designed for
Unicode input and OpenType fonts, and therefore works with XeLaTeX and LuaLaTeX only.

\begin{burmese}
\pkg{myanmar} package က LaTeX စာတမ်းတွေမှာ မြန်မာစာကို Unicode စနစ်နဲ့ 
ရေးလို့ရအောင် လုပ်ပေးထားတဲ့ package လေးပါ။ OpenType font တွေကို 
သုံးထားတဲ့အတွက် XeLaTeX ဒါမှမဟုတ် LuaLaTeX နဲ့ပဲ compile လုပ်လို့ရပါမယ်။
\end{burmese}

The package currently provides:

\begin{itemize}
  \item font selection for common Myanmar fonts;
  \item an inline command for Myanmar text;
  \item a block environment for Myanmar text;
  \item Arabic digit to Myanmar digit conversion;
  \item a basic Burmese localization hook for common document strings.
\end{itemize}

\begin{burmese}
ဒီ package မှာ လူသုံးများတဲ့ မြန်မာ font တွေရွေးတာ၊ စာကြောင်းထဲ မြန်မာစာညှပ်ထည့်တာ၊ 
စာပိုဒ်လိုက် မြန်မာစာရေးတာ၊ ဂဏန်းတွေကို မြန်မာဂဏန်းပြောင်းတာနဲ့ အခြေခံ Burmese localization တွေ ပါဝင်ပါတယ်။
\end{burmese}

This manual describes the public interface implemented by
\texttt{myanmar-script.sty}.

\begin{burmese}
ဒီလမ်းညွှန်က \texttt{myanmar-script.sty} ထဲမှာပါတဲ့ အသုံးပြုသူတွေအတွက် command တွေနဲ့ 
option တွေကို ရှင်းပြပေးထားပါတယ်။
\end{burmese}

% ==================================================
\section{Requirements}

The package requires a Unicode-aware LaTeX engine:

\begin{itemize}
  \item XeLaTeX
  \item LuaLaTeX
\end{itemize}

\begin{burmese}
ဒီ package ကို သုံးဖို့ Unicode ကို အထောက်အပံ့ပေးတဲ့ LaTeX engine လိုအပ်ပါတယ်။ 
XeLaTeX ဒါမှမဟုတ် LuaLaTeX ကို သုံးပေးပါ။
\end{burmese}

The package intentionally stops with an error under pdfLaTeX. Myanmar script
typesetting needs Unicode input and OpenType shaping, which are provided through
\pkg{fontspec}.

\begin{burmese}
pdfLaTeX နဲ့ compile လုပ်ရင်တော့ package က error ပြပြီး ရပ်သွားပါလိမ့်မယ်။ 
မြန်မာစာကို သေသေချာချာ ပုံဖော်နိုင်ဖို့ Unicode input နဲ့ OpenType shaping 
လိုအပ်လို့ \pkg{fontspec} ကို သုံးထားတာပါ။
\end{burmese}

The following LaTeX packages are loaded by \pkg{myanmar-script}:

\begin{itemize}
  \item \pkg{expl3}
  \item \pkg{l3keys2e}
  \item \pkg{iftex}
  \item \pkg{xparse}
  \item \pkg{fontspec}
\end{itemize}

At least one supported Myanmar font must be installed on the system. The
package currently supports \texttt{Noto Sans Myanmar}, \texttt{Pyidaungsu},
\texttt{Myanmar Text}, \texttt{Padauk}, and \texttt{Tharlon}.

\begin{burmese}
System ထဲမှာ support ပေးထားတဲ့ မြန်မာ font အနည်းဆုံးတစ်ခုတော့ install လုပ်ထားရပါမယ်။ 
လက်ရှိ package က \texttt{Noto Sans Myanmar}, \texttt{Pyidaungsu}, 
\texttt{Myanmar Text}, \texttt{Padauk}, နဲ့ \texttt{Tharlon} font တွေကို 
အထောက်အပံ့ပေးထားပါတယ်။
\end{burmese}

% ==================================================
\section{Installation}

The documented package source is \texttt{myanmar.dtx}. The installer is
\texttt{myanmar.ins}. To generate \texttt{myanmar-script.sty}, run:

\begin{Verbatim}
latex myanmar.ins
\end{Verbatim}

\begin{burmese}
Package ရဲ့ source file က \texttt{myanmar.dtx} ဖြစ်ပြီး installer file ကတော့ 
\texttt{myanmar.ins} ပါ။ \texttt{myanmar-script.sty} ကို ထုတ်ယူဖို့ အပေါ်က 
command ကို run ပေးပါ။
\end{burmese}

For local use, place the generated \texttt{myanmar-script.sty} beside your document or
compile from the package root directory.

\begin{burmese}
Local document မှာ သုံးချင်ရင်တော့ ထွက်လာတဲ့ \texttt{myanmar-script.sty} ကို 
ကိုယ့်ရဲ့ \texttt{.tex} file နဲ့ folder အတူတူ ထားပေးပါ။ ဒါမှမဟုတ်လည်း package ရဲ့ root 
directory ကနေ compile လုပ်လို့ရပါတယ်။
\end{burmese}

For a user-level installation, place \texttt{myanmar-script.sty} in a local TeX tree,
for example:

\begin{Verbatim}
texmf/tex/latex/myanmar-script/myanmar-script.sty
\end{Verbatim}

\begin{burmese}
User-level installation အတွက်ဆိုရင်တော့ \texttt{myanmar-script.sty} ကို local TeX tree 
ထဲ ထည့်လို့ရပါတယ်။ TeX distribution အပေါ်မူတည်ပြီး filename database ကို refresh 
လုပ်ပေးဖို့တော့ လိုနိုင်ပါတယ်။
\end{burmese}

% ==================================================
\section{Loading the Package}

Load the package in the preamble:

\begin{Verbatim}
\usepackage{myanmar-script}
\end{Verbatim}

You can also select a font and language:

\begin{Verbatim}
\usepackage[font=noto, language=burmese]{myanmar-script}
\end{Verbatim}

\begin{burmese}
Package ကို document ရဲ့ preamble မှာ load လုပ်ပေးပါ။ Font နဲ့ language ကိုလည်း 
option အနေနဲ့ ရွေးလို့ရပါတယ်။
\end{burmese}

Compile the document with XeLaTeX or LuaLaTeX:

\begin{Verbatim}
xelatex document.tex
lualatex document.tex
\end{Verbatim}

\begin{burmese}
Compile လုပ်တဲ့အခါ XeLaTeX ဒါမှမဟုတ် LuaLaTeX ကိုပဲ သုံးပေးပါ။
\end{burmese}

% ==================================================
\section{Package Options}

The package uses key-value options.

\begin{burmese}
Package option တွေကို key-value ပုံစံနဲ့ သတ်မှတ်ထားပါတယ်။
\end{burmese}

\subsection{\opt{font}}

The \opt{font} option selects the Myanmar font family.

\begin{burmese}
\opt{font} option ကတော့ သုံးချင်တဲ့ မြန်မာ font family ကို ရွေးပေးတာပါ။
\end{burmese}

\begin{center}
\begin{tabular}{ll}
\toprule
Option & Font name \\
\midrule
\opt{font=noto} & \texttt{Noto Sans Myanmar} \\
\opt{font=pyidaungsu} & \texttt{Pyidaungsu} \\
\opt{font=myanmartext} & \texttt{Myanmar Text} \\
\opt{font=myanmar-text} & \texttt{Myanmar Text} alias \\
\opt{font=padauk} & \texttt{Padauk} \\
\opt{font=tharlon} & \texttt{Tharlon} \\
\bottomrule
\end{tabular}
\end{center}

The default is:

\begin{Verbatim}
font=noto
\end{Verbatim}

\begin{burmese}
Default font option ကတော့ \opt{font=noto} ပါ။
\end{burmese}

\subsection{\opt{language}}

The \opt{language} option records the target Myanmar-script language variant.

\begin{burmese}
\opt{language} option ကတော့ သုံးချင်တဲ့ မြန်မာ script language variant ကို 
သတ်မှတ်ပေးတာပါ။
\end{burmese}

\begin{center}
\begin{tabular}{ll}
\toprule
Option & Language \\
\midrule
\opt{language=burmese} & Burmese \\
\bottomrule
\end{tabular}
\end{center}

At present, only the `burmese` value has any implemented localization hooks; other values are not implemented in this release.

The default is:

\begin{Verbatim}
language=burmese
\end{Verbatim}

At present, only \opt{language=burmese} has localization behavior. The other
language values are accepted for future language-specific extensions.

\begin{burmese}
လောလောဆယ်တော့ \opt{language=burmese} အတွက်ပဲ localization အလုပ်လုပ်ဦးမှာပါ။ 
ကျန်တဲ့ language option တွေကတော့ နောက်ပိုင်း တိုးချဲ့လာရင် သုံးလို့ရအောင် 
လက်ခံပေးထားတာပါ။
\end{burmese}

% ==================================================
\section{User Commands}

\subsection{\cmd{textmyanmar}}

Use \cmd{textmyanmar} for short inline Myanmar text:

\begin{Verbatim}
This is inline Myanmar: \textmyanmar{<Myanmar text>}.
\end{Verbatim}

\begin{burmese}
စာကြောင်းထဲမှာ မြန်မာစာ တိုတိုလေးတွေရေးဖို့ \cmd{textmyanmar} ကို သုံးပေးပါ။
\end{burmese}

Rendered example:

\begin{quote}
This is inline Myanmar: \textmyanmar{မြန်မာစာ}.
\end{quote}

\subsection{\texttt{myanmar} environment}

Use the \texttt{myanmar} environment for paragraph or block text:

\begin{Verbatim}
\begin{myanmar}
<Myanmar paragraph text>
\end{myanmar}
\end{Verbatim}

\begin{burmese}
မြန်မာစာပိုဒ်လိုက် ဒါမှမဟုတ် block text အနေနဲ့ရေးဖို့ \texttt{myanmar} environment 
ကို သုံးပေးပါ။
\end{burmese}

Rendered example:

\begin{quote}
\begin{myanmar}
မြန်မာစာသည် မြန်မာနိုင်ငံ၏ ရုံးသုံးဘာသာစကားဖြစ်သည်။
\end{myanmar}
\end{quote}

\subsection{\cmd{myanmarnumber}}

Use \cmd{myanmarnumber} to convert Arabic digits to Myanmar digits:

\begin{Verbatim}
\myanmarnumber{2026}
\end{Verbatim}

\begin{burmese}
အင်္ဂလိပ်ဂဏန်း (Arabic digit) တွေကို မြန်မာဂဏန်း ပြောင်းချင်ရင် \cmd{myanmarnumber} ကို 
သုံးပေးပါ။
\end{burmese}

Rendered example:

\begin{quote}
\myanmarnumber{2026}
\end{quote}

The current implementation performs direct replacement of the characters
\texttt{0} through \texttt{9}. It is not a full number formatter.

\begin{burmese}
လောလောဆယ် လုပ်ပေးထားတာက \texttt{0} ကနေ \texttt{9} အထိ ဂဏန်းတွေကို 
တိုက်ရိုက် အစားထိုးပေးရုံပဲဖြစ်ပြီး full number formatter တော့ မဟုတ်သေးပါဘူး။
\end{burmese}

% ==================================================
\section{Font Setup}

Internally, the package defines:

\begin{Verbatim}
\newfontfamily\myanmarfont{<selected font>}[
  Script=Myanmar,
  Renderer=HarfBuzz,
  Scale=MatchLowercase
]
\end{Verbatim}

The \texttt{Renderer=HarfBuzz} option is useful for LuaLaTeX shaping. XeLaTeX
may warn that this renderer option is LuaTeX-only; that warning is expected and
does not normally prevent output.

\begin{burmese}
Package ထဲမှာတော့ ကိုယ်ရွေးလိုက်တဲ့ font ကို \cmd{myanmarfont} အနေနဲ့ 
သတ်မှတ်ပေးထားပါတယ်။ \texttt{Renderer=HarfBuzz} က LuaLaTeX မှာ shaping 
လုပ်ဖို့အတွက် အသုံးဝင်ပါတယ်။ XeLaTeX သုံးရင်တော့ ဒီ option က LuaTeX အတွက်ပဲဆိုပြီး 
warning တက်လာနိုင်ပေမယ့် ပုံမှန်အားဖြင့်တော့ output ထုတ်တာကို ပြဿနာမပေးပါဘူး။
\end{burmese}

% ==================================================
\section{Localization}

At the beginning of the document, the package runs its localization hook. When
\opt{language=burmese} is selected, it redefines several common document strings
with Myanmar text.

\begin{burmese}
Document စတဲ့အချိန်မှာ package က localization hook ကို run ပါလိမ့်မယ်။ 
\opt{language=burmese} ကို ရွေးထားရင် document ထဲက စာသားတချို့ကို 
မြန်မာလို အလိုအလျောက် ပြောင်းပေးသွားမှာပါ။
\end{burmese}

The current hook covers:

\begin{itemize}
  \item \cmd{contentsname}
  \item \cmd{chaptername}
  \item \cmd{figurename}
  \item \cmd{tablename}
  \item \cmd{bibname}
  \item \cmd{appendixname}
  \item \cmd{indexname}
\end{itemize}

\begin{burmese}
ဒီသတ်မှတ်ချက်တွေက အခြေခံအဆင့်လောက်ပဲ ရှိပါသေးတယ်။ နောက်ပိုင်း version တွေမှာတော့ 
ဗမာ၊ ရှမ်း၊ မွန်နဲ့ စကောကရင် တွေအတွက် ပိုပြီးပြည့်စုံတဲ့ language support တွေ 
ထည့်ပေးလာနိုင်ပါတယ်။
\end{burmese}

% ==================================================
\section{Complete Example}

\begin{Verbatim}
% !TeX program = xelatex
\documentclass{article}

\usepackage[font=pyidaungsu, language=burmese]{myanmar}

\begin{document}

\section{English Header}
This is standard English text.

\section{\textmyanmar{<Myanmar heading>}}
\begin{myanmar}
<Myanmar paragraph text>
\end{myanmar}

Here is a number conversion: \myanmarnumber{2026}

\end{document}
\end{Verbatim}

\begin{burmese}
အပေါ်က sample မှာ \opt{font=pyidaungsu} ကို သုံးပြထားပါတယ်။ ကိုယ့် system မှာ 
Pyidaungsu မရှိဘူးဆိုရင် \opt{font=noto}, \opt{font=myanmartext}, 
\opt{font=padauk}, ဒါမှမဟုတ် \opt{font=tharlon} ကို ပြောင်းသုံးလို့ရပါတယ်။
\end{burmese}

% ==================================================
\section{Building This Manual}

From the repository root, run:

\begin{Verbatim}
xelatex docs/manual.tex
\end{Verbatim}

or:

\begin{Verbatim}
lualatex docs/manual.tex
\end{Verbatim}

If \texttt{latexmkrc} is present, this should also work:

\begin{Verbatim}
latexmk docs/manual.tex
\end{Verbatim}

\begin{burmese}
ဒီ manual ကို build လုပ်ချင်ရင် repository ရဲ့ root ကနေ အပေါ်က command တွေကို 
run ပေးပါ။ \texttt{latexmkrc} ဖိုင်ရှိရင်တော့ \texttt{latexmk} က XeLaTeX သုံးဖို့ 
အသင့်ပြင်ပေးထားနိုင်ပါတယ်။
\end{burmese}

% ==================================================
\section{Troubleshooting}

\subsection{The package says pdfLaTeX is not supported}

Compile with XeLaTeX or LuaLaTeX instead:

\begin{Verbatim}
xelatex document.tex
\end{Verbatim}

\begin{burmese}
pdfLaTeX ကို မသုံးပါနဲ့။ XeLaTeX ဒါမှမဟုတ် LuaLaTeX နဲ့ပဲ compile လုပ်ပေးပါ။
\end{burmese}

\subsection{A Myanmar font cannot be found}

Install the selected font or choose another supported font option:

\begin{Verbatim}
\usepackage[font=pyidaungsu]{myanmar}
\end{Verbatim}

\begin{burmese}
ရွေးလိုက်တဲ့ font ကို ရှာမတွေ့ရင် အဲဒီ font ကို အရင် install လုပ်ပေးပါ။ 
ဒါမှမဟုတ်လည်း package က အထောက်အပံ့ပေးထားတဲ့ တခြား font option တစ်ခုခုကို 
ရွေးသုံးလို့ရပါတယ်။
\end{burmese}

\subsection{Myanmar text appears as empty boxes}

The active font may not contain Myanmar glyphs. Use one of the supported
Myanmar fonts and wrap Myanmar text in \cmd{textmyanmar} or the
\texttt{myanmar} environment.

\begin{burmese}
စာလုံးတွေက အတုံးလေး (box) တွေလို ဖြစ်နေရင် ကိုယ်သုံးနေတဲ့ font မှာ Myanmar glyph 
မပါလို့ဖြစ်နိုင်ပါတယ်။ Package က ထောက်ပံ့ပေးထားတဲ့ မြန်မာ font တစ်ခုခုကို သုံးပြီး 
မြန်မာစာတွေကို \cmd{textmyanmar} ဒါမှမဟုတ် \texttt{myanmar} environment ထဲ 
ထည့်ရေးပေးပါ။
\end{burmese}

\subsection{Shaping looks incorrect}

Use XeLaTeX or LuaLaTeX with a proper OpenType Myanmar font. For LuaLaTeX,
HarfBuzz rendering is requested by the package.

\begin{burmese}
စာလုံးပုံစံတွေ အချိတ်အဆက် မမှန်ရင် OpenType Myanmar font သုံးပြီး XeLaTeX 
ဒါမှမဟုတ် LuaLaTeX နဲ့ compile လုပ်ကြည့်ပါ။ LuaLaTeX သုံးမယ်ဆိုရင် package က 
HarfBuzz rendering ကို တောင်းဆိုထားပါတယ်။
\end{burmese}

% ==================================================
\section{Current Limitations}

\begin{itemize}
  \item Only five font choices are currently predefined, plus the
    \opt{font=myanmar-text} alias.
  \item Language-specific behavior is currently implemented only for Burmese.
  \item Number conversion is simple digit replacement.
  \item The package does not provide Myanmar line-breaking or hyphenation data.
  \item The package does not yet provide separate language mapping files.
\end{itemize}

\begin{burmese}
လက်ရှိ version မှာ font choice ငါးခုနဲ့ \opt{font=myanmar-text} alias ကိုပဲ 
ထည့်ပေးထားပါသေးတယ်။ Language-specific behavior အနေနဲ့လည်း ဗမာစာ (Burmese) အတွက်ပဲ 
အခြေခံအဆင့် ပါပါသေးတယ်။ ဂဏန်းပြောင်းတဲ့နေရာမှာလည်း digit ကို အစားထိုးပေးရုံပဲဖြစ်ပြီး 
line-breaking တွေ၊ hyphenation data တွေတော့ မပါသေးပါဘူး။
\end{burmese}

% ==================================================
\section{License}

This repository is distributed under the LaTeX Project Public License (LPPL) --- Version 1.3c (2008-05-04). See the \texttt{LICENSE}
file for the full license text.

\begin{burmese}
ဤ repository သည် LaTeX Project Public License (LPPL) ဗားရှင်း 1.3c (2008-05-04) အရ ဖြန့်ဝေထားသည်။ နောက်ထပ် license အပြည့်အစုံကို \texttt{LICENSE} ဖိုင်ထဲမှ ဖတ်ရှုနိုင်ပါသည်။
\end{burmese}

\end{document}
